\section{The all-pilot palette and the local zero-socket test}
\label{sec:actual-palette-warning}

This section isolates the warning input used by the opacity construction.
There are two different estimates here. The palette estimate is a global,
exponentially weighted row-and-column bound for the actual screened Gaussian
matrix. The socket conclusion is deterministic: on a prescribed local
small-noise event, every factor in the local mask product equals one. A
comparison for one posed reader does not imply the first estimate, and a
Gaussian probability estimate does not by itself imply the second.

For $x>1$ write $\operatorname{radius}(x)=\lceil(\log x)^2\rceil$, the
source radius $\operatorname{radiusB}(x)$.

Write $\mathcal I$ for the active link-colour coordinates of a fixed forest
and $\mathcal R$ for the retained response-row coordinates. A move matrix
$M=(M_{ri})_{r\in\mathcal R,i\in\mathcal I}$ maps active-coordinate
vectors to response-row vectors; $a(r)$ and $a(i)$ denote their anchor
sites. For a rectangular matrix define its exponential row-and-column price
at rate $\sigma\geq0$ by
\[
 \mathfrak p_\sigma(M)=\max\left\{
   \sup_{r\in\mathcal R}\sum_{i\in\mathcal I}
       e^{\sigma d(a(r),a(i))}|M_{ri}|,
   \sup_{i\in\mathcal I}\sum_{r\in\mathcal R}
       e^{\sigma d(a(r),a(i))}|M_{ri}|\right\}.
\]
The bounds below are uniform in the torus size, but their constants may
depend on the fixed patch, tube, and interaction range. Let $Q$ be the Gram
matrix of the actual row family, let $W$ be the panel projection, and put
\[
 C=\operatorname{clouded}(W,Q),\qquad N=\operatorname{per}(WQ).
\]
Here $s\in\mathbb R^{\mathcal R}$ is the source row,
$b=M^{\mathsf T}s\in\mathbb R^{\mathcal I}$ its transported row and
$m=-N^{\mathsf T}b$ the selected mean. The corresponding direct and
auxiliary rows are
\[
 r=Mm+s,\qquad
z^{\rm aux}=-(\operatorname{coupling}(\beta))^{-1}
       Q\bigl(Q^{-1}-N^{\mathsf T}\bigr)b.
\]
These are the finite-dimensional identities behind the estimates; the
transpose is the ordinary real transpose. The windowed covariance $C$ and
its inverse comparison are taken on the same retained coordinate set.

\begin{proposition}[Weighted response implies a uniform palette bound]
\label{prop:actual-palette-response}
Suppose that, for all sufficiently large $x$, uniformly in the period,
forest, tube, choices and configuration under consideration, the following
prices hold:
\begin{enumerate}
\item $Q$ and $C$ are positive definite and
      $\mathfrak p_0(Q),\mathfrak p_0(C),\mathfrak p_0(N)\leq K$;
\item $\mathfrak p_0(M)\leq L$;
\item $|s_r|\leq Z\,\varphi(x)$ at every source coordinate $r\in\mathcal R$, where
      $K,L,Z>0$, $\varphi(x)\geq0$, and $\varphi(x)\to0$; the local
      width satisfies $0\leq d\leq r\varphi(x)$ for a fixed $r\geq0$;
\item the window defect obeys
\[
 \mathfrak p_0\bigl(N-Q^{-1}\bigr)
       \leq u\exp\{-v\,\operatorname{radius}(x)/60\}
\]
for fixed $u,v>0$;
\item $0<\beta\leq4x$ and $x\geq4$. Since
      $g=\operatorname{coupling}(\beta)=\beta^{-1/2}$, we have
      $g^{-1}=\sqrt\beta\leq2\sqrt x\leq x$. The local door product
      is a product of broad factors and obeys this implication: if it is
      nonzero, then $|m_i+g\zeta_i|\leq1$ whenever $|m_i|\leq1$ and
      $2d\leq1$, where $g=\operatorname{coupling}(\beta)$ and $\zeta$
      is the noise vector.
\end{enumerate}
Then there is a fixed $\ell\geq1$ such that, eventually in $x$, the
following four bounds hold for every configuration row in this family:
\[
 |m_i|\leq1,\qquad |r_i|\leq1,\qquad |z_i^{\rm aux}|\leq x^{-4},
 \qquad \text{door product}\ne0\Longrightarrow
 |\zeta_i|\leq x^2.
\]
Moreover the screened palette matrix $C$ satisfies
\[
 \mathfrak p_0(C)\leq\ell.
\]
\end{proposition}
\begin{proof}
For any rectangular matrix $A$ with $\mathfrak p_0(A)\leq c$ and any vector
$u$ on its column set with $\|u\|_\infty\leq b_0$, the triangle inequality gives
\[
 |(Au)_i|\leq\sum_j|A_{ij}|\,|u_j|
             \leq c b_0.
\]
The move-map price first gives
$|b_i|=|(M^{\mathsf T}s)_i|\leq LZ\varphi(x)$ by its column bound.
Apply the row-and-column estimate to $N^{\mathsf T}b$. The column bound
for $N$ is the row bound for its transpose, so
$|m_i|\leq KLZ\varphi(x)$. Applying the same calculation to $Mm+s$ gives
\[
 |r_i|\leq (L^2KZ+Z)\varphi(x).
\]
For the auxiliary row, factor the displayed identity exactly as written:
\[
 |z_i^{\rm aux}|\leq x\,K u LZ\varphi(x)
       e^{-v\operatorname{radius}(x)/60}.
\]
Because $\operatorname{radius}(x)$ grows quadratically in $\log x$,
the exponential factor is eventually at most $x^{-5}$. The right side is
then at most $K u LZ\varphi(x)x^{-4}$. Since $K,u,L,Z$ are fixed and
$\varphi(x)\to0$, its coefficient is eventually at most one; hence
$|z_i^{\rm aux}|\leq x^{-4}$. The first two row bounds also tend to zero,
hence are eventually at most one.
The width condition and $\varphi(x)\to0$ ensure $2d\leq1$ eventually.
If the door product is nonzero, its broad-span implication gives
$|m_i+g\zeta_i|\leq1$. Thus $|g\zeta_i|\leq2$, and the coupling-price
hypothesis gives
\[
 |\zeta_i|\leq2g^{-1}\leq2x\leq x^2
\]
for $x\geq2$. Finally, the weighted price of $C$ is already at most $K$;
take $\ell=\max(1,K)$. This proves all conclusions with one threshold,
independent of the period and of the subsequently chosen configuration.
\end{proof}

The four conclusions, together with $\mathfrak p_0(C)\leq\ell$, are the
complete response receipt needed here: they give a positive weighted
palette price, a mean bound, a direct-row bound, an auxiliary-row bound
of order $x^{-4}$, and the implication from a nonzero door product to the
noise bound $x^2$. Section~\ref{sec:actual-global-row-prices} supplies
the distinct global matrix, source-row, and move-row estimates used in
this calculation.

In the selected application, Section~\ref{sec:actual-global-row-prices}
supplies the matrix price and gap for $Q$, its panel restriction $C$, the
windowed matrix $N$, and the exponentially small price for $N-Q^{-1}$.
It also supplies the source-row bound $Z\operatorname{fiery}(x)$ and the
move-row price $L$. The calculation above is included because it shows
exactly how these different receipts combine, with no factor depending on
the number of sites. The local covariance comparison of
Section~\ref{sec:actual-reader-local-gaussian-prices} is a distinct input;
its entrywise estimate cannot replace these global weighted bounds.

\subsection{Why the palette estimate is all-matched}

For a pilot $p$, let $\Pi_p$ be its endpoint palette, namely the normalized
Gaussian matrix on the forest coordinates after retaining the panel
addresses. A palette has a $\sigma$-price $C$ when its two weighted sums
above are at most $C$. The endpoint convention identifies $\Pi_p$ exactly
with the clouded panel Gram matrix $C$; it is not an approximate
identification and involves no limiting field.

For each matched pilot, specialize the uniform response estimate to that
pilot's store, forest, configuration and tag. The endpoint palette then
satisfies
\[
  \mathfrak p_0(\Pi_p)\leq\ell\qquad\text{for every matched }p.
\]
Because the matrix prices were obtained uniformly before selecting the
pilot, the same constant works for the entire matched family. This
universal quantifier comes from all-configuration screened-kernel bounds,
not from testing a dense set of readers or passing from one posed pilot by
continuity.

\subsection{The local zero-socket implication}

Fix a matched pilot $p$ and a mask address $i$, and put
$q_i=\operatorname{toes}(i)$ and $B=\operatorname{radius}(x)$. Write
$z_p$ for its local fluctuation vector and $g=\operatorname{coupling}(\beta)>0$.
For a site $o$, define $\mathsf{Toned}(R,o,z,h)$ to mean
$|z_j|\leq h$ for each active coordinate $j$ whose anchor is at distance
at most $R$ from $o$. The local small-noise condition is
\[
 \mathsf{Toned}(2B,q_i,z_p,\varepsilon\,t/g),
 \qquad t=\operatorname{fiery}(x).
\]
By definition, the
socket is zero when the support foot meets the store's radius-$300B$ dome.
It remains to treat the branch in which the foot and that dome are
disjoint. The foot contains its toe $q_i$. Since the radius-$50B$ dome is contained
in the radius-$300B$ dome, disjointness immediately gives
$q_i\notin A.dome(50B)$.

The beacon receipt is evaluated on the actual selected configuration,
store, tag and mean-adjusted field. At every site $y$ outside the
radius-$40B$ store dome, it gives active tape coordinates of norm below
the chart radius, fine-face toll at most $t/4$, coarse-bond wire at most
$Wt$, and both normalized likelihoods equal to one, provided
$\mathsf{Toned}(B,y,z_p,\varepsilon t/g)$ holds.
The toe is outside the radius-$50B$ dome. Every site within distance $B$ of
the toe is therefore outside the radius-$40B$ dome. Restricting the
radius-$2B$ toned bound gives the required radius-$B$ bound at these sites.
The face, bond-wire and frame parts of the beacon receipt consequently
give the local outpost condition.

The port estimate turns this local condition into the door bound. Let
$N_{\mathrm{port}}$ be the number of plaquettes in the fixed macrobox for
the patch (the macrobox site count times the number of plaquette
orientations). The port inequality sums at most $N_{\mathrm{port}}$ face
tolls, and the braid inequality
bounds the connector by $Wt$. Thus the chord of the relative link at any
active row is at most $(C_0N_{\mathrm{port}}+W)t$, where $C_0$ is the fixed port constant.
On the chart ball, the logarithm has norm at most $K_0$ times this chord.
The beacon places the actual relative link in the injectivity ball, so its
logarithm is the actual tape coordinate. Therefore the outpost condition and
the distant-door estimate give
\[
 \|\ell_{p,j}\|\leq C_*t,
 \qquad C_*=K_0(C_0N_{\mathrm{port}}+W),
\]
for each active door row. An inactive door factor equals one by its
definition. The selected width has $C_*t\leq\operatorname{arc}$, so the
broad factor equals one for active rows too. The integer $N_{\mathrm{port}}$ is fixed by the patch and does not grow with the torus.

The mask has five coordinate types: fine-face share, coarse-bond wire
share, site likelihood, bond likelihood, and door. The first two equal one
because their nonnegative tolls are below their full-share thresholds; the
likelihoods equal one by the beacon receipt; and the door equals one by the
preceding argument. Thus the local ask is one, so the socket (ask minus
one) is zero. Every estimate is on the same matched pilot; there is no
perturbation or replacement of its store or tag.

\begin{proposition}[Exact assembly of the palette condition]
\label{prop:actual-palette-condition}
Assume $\varepsilon\geq0$ and $\beta>0$, and the uniform weighted response prices of
Proposition~\ref{prop:actual-palette-response}; these give a common price
bound $\ell\geq1$. Let $\mathcal P$ be the set of matched pilots. Assume
the selected field and forest satisfy the same-history response conditions
of Section~\ref{sec:actual-global-row-prices}, and that the following
uniform local receipts hold for every $p\in\mathcal P$ and mask address:
the disjoint-foot branch geometry, beacon estimates, port estimate, and
distant-door estimate just stated. These receipts are on the same field,
forest, tag and store, with thresholds uniform over the matched family.
Then for
\[
 q=\varepsilon\,\operatorname{fiery}(x)/
                   \operatorname{coupling}(\beta)
\]
the endpoint palette has a uniform positive Pours price for every matched
pilot, and its socket vanishes whenever the specified $\mathsf{Toned}(2B,q_i,z_p,\varepsilon t/g)$ condition holds.
These are the two components of the all-matched palette condition used by the opacity step.
\end{proposition}
\begin{proof}
The weighted matrix estimate is uniform over the matched family, so its
constant is the common Pours price $\ell$. For each matched pilot and each
mask address, the local argument above gives zero socket
under its toned hypothesis. Since neither conclusion changes the pilot,
the quantified statements hold for the full matched family.
\end{proof}

\subsection{The selected all-pilot package}

The geometric and local receipts are selected together. For a fixed
patch, a permitted chart ratio, a positive mask width $W$, and the reserve
and wedge hypotheses of the sector construction, there is a lower scale
$L_0$ and a beacon parameter $\varepsilon>0$. For each fixed chosen
width $\lambda\geq L_0$, the chart and size conditions, the annular
response at width $\lambda$, the local beacon, and the distant-door
estimate all hold eventually in $x$. Their cutoff is uniform over legal
stores, forests, tags, configurations and mask addresses, hence over
matched pilots. The port/log-chart and distant-door route gives the fixed
constant $C_*$ above; choose $L_0\geq C_*$, so
$C_*t\leq\lambda t=\operatorname{arc}$. The annular estimate is a
separate component of the selected package.

Here is the mean-switch argument that makes the beacon a statement about
the actual sewn field. The centered affine background (CAB) calibration is
an auxiliary calculation;
its matrices must not be identified with the actual-reader matrices
$Q,C,N$ above. Fix one selected store, forest, tag, field $V$, and a site
$o$ outside the radius-$30B$ store dome, where $B=\operatorname{radiusB}(x)$.
Write $\varepsilon_0=2\varepsilon$, where $\varepsilon$ is the final
all-matched beacon fee and $\varepsilon_0$ is the CAB plateau threshold.
The based chart construction supplies a gauge and a chart vector $b$ with
$\|b\|\leq\operatorname{clearance}(x)/2$, agreeing with that gauge on the
radius-$20B$ room. Write $v$ for the corresponding local affine direction.
For the auxiliary finite calibration, write $H$ for its sheet,
$P_{\rm CAB}$ for its within-sheet map,
$E=HP_{\rm CAB}-I$ for its error, and $s_{\rm CAB}$ for its source.
The exact finite-dimensional subtraction identity is
\begin{equation}\label{eq:actual-palette-cab-subtraction}
 \mu_{\rm CAB}-v
   =-P_{\rm CAB}^{\mathsf T}(Hv+s_{\rm CAB})+E^{\mathsf T}v.
\end{equation}
It follows by expanding $E=HP_{\rm CAB}-I$ and the definitions of the centered
calibration mean. To state the paid rates without source shorthand, put
\[
 \rho_{r,k}(x)=\operatorname{radiusB}(x)^k
                    \,\operatorname{fiery}(x)^r.
\]
The bought central residual and error-sheet estimates are, respectively,
the rates $\rho_{2,25}$ and $\rho_{2,26}$;
the selected scale schedule includes $\rho_{2,27}(x)\leq x^{-7/8}$
eventually. In the first term of (\ref{eq:actual-palette-cab-subtraction}),
the central residual bound $\rho_{2,25}$ is multiplied by the fixed
within-sheet row price; one additional factor of the radius is absorbed
into the next rate, giving $\rho_{2,26}$. The second term is the
error-sheet tail, also at that rate. The reserve's allowance growth then
gives, coordinatewise,
\begin{equation}\label{eq:actual-palette-cab-velocity}
 |(\mu_{\rm CAB}-v)_j|\leq \rho_{2,27}(x)\leq x^{-7/8}.
\end{equation}
These are the two matrix terms in (\ref{eq:actual-palette-cab-subtraction});
no equality between the auxiliary calibration law and the actual palette
law is used.

The chart comparison of the actual portrait with this same affine direction
has a second, separately paid error. The finite chart estimate and the
identity between the local face read and its short-row read give a
per-color error at most $x^{-7/8}$; summing over the fixed color
coordinates and combining with (\ref{eq:actual-palette-cab-velocity}) gives
\begin{equation}\label{eq:actual-palette-mean-switch}
 \|D'_j-\mu_{{\rm CAB},j}\|
       \leq 2\dim(G_{\rm color})x^{-7/8}\leq x^{-3/4}
\end{equation}
for all sufficiently large $x$. Here $D'$ is the portrait of the original
selected field, in the local chart, and the last inequality is just a
fixed-dimension power comparison. If the local fluctuation is $z$ and
$g=\operatorname{coupling}(\beta)>0$, the switched active coordinate is
exactly
\[
   (D'+gz-\mu_{\rm CAB})_j.
\]
For any point in the radius-$B$ neighborhood of $o$,
(\ref{eq:actual-palette-mean-switch}) and the Toned bound
$|z_j|\leq\varepsilon_0t/(2g)$ imply
\[
 |(D'+gz-\mu_{\rm CAB})_j|
      \leq x^{-3/4}+\varepsilon_0t/2\leq\varepsilon_0t
\]
eventually; the final step uses the bought threshold
$x^{-3/4}\leq\varepsilon_0t/2$. Thus the shifted affine field is Awake
with threshold $\varepsilon_0t$ on that neighborhood.

The charged plateau estimate is applied to this auxiliary affine field.
Write $W_{\rm wire}=W$, $a_* = \sqrt{\mathrm{direct}}$, and
$b_* = \sqrt{\mathrm{penalty}}$. The reserve and wedge purchases give
$a_*>0$, $b_*>0$, and the uniform row price satisfies $W_{\rm row}>0$;
these are the positive quadratic and penalty floors in the selected
calibration.
\begin{equation}\label{eq:actual-palette-cab-wedge}
 \frac{C}{L^2}\leq\frac{a_*}{64},\qquad
 C+\frac{a_*}{64}\leq\frac{W_{\rm wire}b_*}{4}.
\end{equation}
Let $W_{\rm row}$ be the bought uniform row price for the selected CAB
system and let $\omega_{\rm CAB}>0$ be its fixed quadratic-response weight.
Set $\delta=a_*/64$ and choose the plateau threshold
\begin{equation}\label{eq:actual-palette-cab-target-fee}
 \varepsilon_0=\min\left\{1,
   \frac{\delta\sqrt{\omega_{\rm CAB}}}
   {4(1+\sqrt{\dim(G_{\rm color})})W_{\rm row}}\right\}.
\end{equation}
The local target estimate gives target size at most $\delta t$ under this
shift. Here is its payment. The action of the CAB mean is bounded by the
rate $\rho_{2,29}$, which is eventually at most
$W_{\rm row}\varepsilon_0t$; the sparse shifted action has the same bound.
Thus the full shifted action is at most
$2W_{\rm row}\varepsilon_0t$. In the target expansion, the linear term is
at most $\delta t/2$ by the choice in (\ref{eq:actual-palette-cab-target-fee}),
and the nonlinear term is at most $\rho_{2,40}\leq\delta t/2$ by the paid
rate comparison. Their sum is at most $\delta t$. Decreasing a permitted
shift threshold preserves these inequalities, so the minimum with one in
(\ref{eq:actual-palette-cab-target-fee}) is harmless.

The first inequality in (\ref{eq:actual-palette-cab-wedge}) bounds the
unshifted face response by $a_*t/64$. The paid target response changes it
by at most $a_*t/64$, so the resulting row has norm at most $a_*t/32$.
On the purchased chart, the source curvature lower bound therefore gives
face toll below $t/4$. For a bond row the unshifted term is at most $Ct$
and the response term at most $a_*t/64$. Its normalized read is $b_*$
times its curvature norm, so the second inequality in
(\ref{eq:actual-palette-cab-wedge}) gives curvature at most
$W_{\rm wire}t/4$. Chord is at most twice curvature on this chart; hence
wire is below $W_{\rm wire}t$. The likelihood plateau estimate sets the
site and bond factors equal to one.

For the broad chart, the reserve bounds each CAB-mean color coordinate by
$\rho_{1,12}$. The shift coordinates are at most $\varepsilon_0t$.
Because $\varepsilon_0\leq1$ and $\operatorname{radiusB}(x)\geq1$,
$\varepsilon_0t\leq\rho_{1,13}$ and $\rho_{1,12}\leq\rho_{1,13}$.
The Euclidean color norm of the shifted row is at most
$2\sqrt{\dim(G_{\rm color})}\rho_{1,13}$, which the purchased chart
budget bounds by the chart radius. Thus the single threshold
$\varepsilon_0$ ensures the target, face, wire and chart conclusions for
every allowed calibration and every local shift. The coordinate-to-color
step is the elementary inequality
$\|a\|\leq\sqrt{\dim(G_{\rm color})}\max_j|a_j|$.

Finally the CAB construction is sewn back without changing the physical
field. Let $\mathcal B$ be the auxiliary configuration map built from the
original field $V$ and its frozen-link data, and let $a=D'+gz$ be the desired
active read. Define $\chi(a)$ to be $a-\mu_{\rm CAB}$ on active coordinates
and zero on frozen coordinates. The sewing identity is, link by link,
\begin{equation}\label{eq:actual-palette-cab-shelter}
 \mathcal B(\mu_{\rm CAB}+\chi(a))=V_{\rm sewn}(a),
\end{equation}
where $V_{\rm sewn}(a)$ uses the prescribed tape $a$ on active links and
the original stored values on frozen links. On a trunk both sides are the
identity. On an active link the auxiliary mean cancels and leaves exactly
the prescribed tape $a$. On a frozen link the correction is zero and the
stored link is unchanged. Thus the plateau bounds
apply to the actual sewn configuration. This proves the local beacon
receipt used in the socket argument; its uniformity in all stores,
forests, tags and periods is the uniformity of the chart, reserve and
wedge purchases, before the pilot is chosen.

Now fix one such selected field, with $0<\beta\leq4x$, local
width $0<d\leq u t$, chosen broad-response width $\lambda\geq L_0$,
and arc $\lambda t$. Suppose its budget certificate gives
both the chart compatibility needed for the beacon and the auxiliary
sector comparison, and suppose the all-matched response estimate above
holds. For each matched pilot the actual response receipt gives the
endpoint Pours bound with the common constant $\ell$. The deterministic
argument above gives zero socket under
$\mathsf{Toned}(2B,q_i,z_p,\varepsilon t/\operatorname{coupling}(\beta))$.
The fee $\varepsilon t/\operatorname{coupling}(\beta)$ is nonnegative.
Thus one and the same selected field satisfies both palette clauses for
all matched pilots. This is the exact two-clause output used in the later
opacity estimate; it is not obtained by selecting a different field for
each pilot.

The source's assembly uses these inputs in a fixed order. An enlarged
patch supplies the local chart, while the separate port/log-chart estimate
and the distant-door estimate supply the far-door bound. The annular
response supplies its own fixed-scale response receipt; it is not the source
of the port constant.
Combining these with the all-configuration weighted palette estimate gives
the two hypotheses of the proposition on one selected sector. The
correlated-warning bound of Section~\ref{sec:gaussian-warnings} prices the
probability of violating many such local conditions, while
Section~\ref{sec:warning-safety} proves the deterministic likelihood and
row plateaus on the small-fluctuation event. The annular argument in
Section~\ref{sec:packet-annulus} controls the energy outside the packet
core. None of these three statements alone is the all-matched Pours
estimate; that part remains tied to the global weighted-kernel supplier.
